\documentclass[11pt]{article}

\usepackage[margin=1in]{geometry}
\usepackage[utf8]{inputenc}
\usepackage[T1]{fontenc}
\usepackage{lmodern}
\usepackage{xcolor}
\usepackage{textcomp}
\usepackage{titlesec}
\usepackage{enumitem}
\usepackage{parskip}
\usepackage[colorlinks=true, linkcolor=blue!50!black, citecolor=blue!50!black, urlcolor=blue!50!black]{hyperref}

\usepackage[backend=biber, style=numeric, sorting=none]{biblatex}
\addbibresource{references.bib}

\titleformat{\section}{\Large\bfseries}{\thesection}{0.6em}{}
\titleformat{\subsection}{\large\bfseries}{\thesubsection}{0.6em}{}
\titlespacing*{\subsection}{0pt}{1.4em}{0.4em}
\titlespacing*{\paragraph}{0pt}{1.3em}{0.5em}

\setlist[description]{leftmargin=1.4em, itemsep=0.7em, font=\normalfont\bfseries}

\newcommand{\charname}[1]{\textbf{#1}}
\newcommand{\subcat}[1]{\par\medskip\noindent\textit{#1}\par\smallskip\noindent}

\title{\textbf{A Reader's Companion to \emph{1929}}\\[0.3em]
\large Andrew Ross Sorkin's \emph{1929: Inside the Greatest Crash in Wall Street History \\ and How It Shattered a Nation}}
\author{Prepared by Claude (Sonnet 5), Anthropic \\ \small for \href{mailto:icos.atropa@gmail.com}{a listener of the audiobook narration}}
\date{August 24, 2026}

\begin{document}

\maketitle

\tableofcontents

\section{About}
\label{sec:about}

Andrew Ross Sorkin is a financial journalist and the co-anchor of CNBC's
\emph{Squawk Box}, best known previously for \emph{Too Big to Fail} (2009), his
account of the 2008 financial crisis \autocite{sorkin2025book,prh1929,goodreads1929,
airmail1929}. In his own author's note, Sorkin situates \emph{1929} alongside three
works he calls foundational to the field: John Kenneth Galbraith's \emph{The Great
Crash 1929} (1955), John Brooks's \emph{Once in Golconda} (1969), and Gordon Thomas
and Max Morgan-Witts's \emph{The Day the Bubble Burst} (1979)
\autocite{sorkinPreviewPDF}. He distinguishes his own contribution by more than
eight years of reporting that included first-time access to the Federal Reserve
Bank of New York's confidential 1929 board minutes and a previously unpublished
memoir from a Wall Street insider, and by a deliberately character-driven
approach: the book follows the motivations of, in his words, ``those who helped
set [the crash] in motion,'' rather than retelling the crisis through statistics
and policy alone \autocite{sorkinPreviewPDF,axios1929,realclearmarkets1929}.
Reviewers have described the result as a kind of prequel to \emph{Too Big to
Fail}, reading 1929 and 2008 as two chapters in the same recurring story of
Wall Street risk-taking and its consequences \autocite{airmail1929}.

\section{Cast of Characters}
\label{sec:cast-characters}

The figures below are the major, named individuals the book follows closely across
its narrative, drawn from publisher and reviewer descriptions of the book's
principal cast \autocite{supersummary1929,supersummaryKeyFigures,booksthatslay1929,
littlerbooks1929}. Wilson and Coolidge, who appear only briefly as outgoing political
context, are omitted at the reader's request.

\begin{description}

\item[Charles E.\ Mitchell] Chairman of \charname{National City Bank} (a
predecessor of Citibank) and known as ``Sunshine Charlie'' for his relentless
public optimism, Mitchell pioneered the mass marketing of stocks to ordinary
Americans and used his bank's resources to keep credit flowing to the market
even when the \charname{Federal Reserve} wanted to restrain it
\autocite{airmail1929,supersummary1929}. His running conflict with the
\charname{Federal Reserve}'s leadership is one of the book's central threads
connecting Wall Street to Washington.

\item[Thomas W.\ Lamont] The senior partner at \charname{J.P.\ Morgan \& Co.}\
and the most influential private banker of the era, Lamont acted as an informal
statesman for Wall Street, coordinating with fellow bankers and government
officials to manage the market's mood \autocite{supersummary1929,
dailyeconomy1929}. He is a connective figure between the \charname{House of
Morgan}, the \charname{New York Stock Exchange}, and the \charname{Hoover}
administration.

\item[J.P.\ Morgan Jr.] Head of the \charname{House of Morgan} and heir to his
father's banking dynasty, Morgan represents the old-guard establishment of
American finance, whose firm's private influence over markets and government
rivaled that of the \charname{Federal Reserve} itself
\autocite{supersummary1929}. He works closely with \charname{Lamont} throughout
the book.

\item[Richard Whitney] A vice president (and later president) of the
\charname{New York Stock Exchange}, Whitney became the public face of Wall
Street's efforts to project confidence on the trading floor during the
market's most volatile days, acting on behalf of the leading bankers'
consortium \autocite{supersummary1929,csmonitor1929}. He works in close
coordination with \charname{Lamont} and the Morgan-led banking group.

\item[Jesse Livermore] A legendary independent stock trader famous for making
and losing several fortunes, Livermore operated outside the banking
establishment, trading on his own instincts about market psychology rather than
on inside information or institutional loyalty \autocite{airmail1929,
supersummary1929}. He functions in the narrative as a skeptical outsider
watching the same boom that bankers like \charname{Mitchell} were fueling.

\item[William C.\ Durant] The founder of \charname{General Motors}, Durant
reinvented himself in the late 1920s as one of Wall Street's most aggressive
stock market operators, organizing pools of investors to move prices in
favored stocks \autocite{supersummary1929}. His path connects
\charname{General Motors}, the speculative pool culture of the exchange, and
fellow GM figure \charname{John J.\ Raskob}.

\item[John J.\ Raskob] A \charname{General Motors} and DuPont executive who
also chaired the Democratic National Committee, Raskob was a leading public
evangelist for stock ownership, urging ordinary Americans that ``everybody
ought to be rich'' through the market \autocite{supersummary1929,
littlerbooks1929}. He links the corporate world of \charname{General Motors}
to both Wall Street speculation and national politics.

\item[Michael J.\ Meehan] A \charname{New York Stock Exchange} specialist and
stock promoter known for running trading pools in popular speculative stocks
such as Radio Corporation of America (RCA), Meehan exemplifies the era's
loosely regulated market-manipulation practices \autocite{supersummary1929}.
His activities intersect with the same pool culture that drew in
\charname{Durant} and \charname{Raskob}.

\item[Herbert Hoover] Inaugurated president in March 1929, Hoover inherited a
market boom he privately distrusted and then had to respond to its collapse,
caught between his instincts for voluntary, business-led solutions and mounting
pressure for federal action \autocite{supersummary1929,csmonitor1929}. His
administration's dealings with bankers like \charname{Lamont} and
\charname{Mitchell} are a recurring axis of the book.

\item[Andrew W.\ Mellon] A wealthy financier and industrialist serving as
Treasury Secretary, Mellon represents the fusion of private fortune and public
economic policy that characterized the era, shaping the government's hands-off
approach to the boom \autocite{supersummary1929}. He works alongside
\charname{Hoover} through the crash's early aftermath.

\item[George L.\ Harrison] Governor of the \charname{Federal Reserve Bank of
New York}, Harrison inherited the internal \charname{Federal Reserve} debates
over whether and how to curb speculative lending, and dealt directly with
bankers such as \charname{Mitchell} whose actions the Fed was trying to
restrain \autocite{supersummary1929}. He is the book's throughline for the
\charname{Federal Reserve}'s institutional perspective.

\item[Franklin D.\ Roosevelt] Governor of New York during the crash and elected
president in 1932, Roosevelt becomes central to the book's second half as the
political figure whose administration presided over the era's regulatory
response to Wall Street's excesses \autocite{supersummary1929,libraryjournal1929}.
He connects to Senator \charname{Carter Glass} and investigator
\charname{Ferdinand Pecora} as reform-era figures.

\item[Carter Glass] A U.S.\ Senator and former Treasury Secretary, Glass was a
leading architect of banking reform legislation in the crash's aftermath,
working to redraw the boundary between commercial and investment banking
\autocite{libraryjournal1929}. His legislative efforts connect Washington's
political response to the practices of bankers like \charname{Mitchell} and
\charname{Lamont} profiled earlier in the book.

\item[Ferdinand Pecora] Chief counsel to the \charname{U.S.\ Senate Committee
on Banking and Currency}, Pecora led the high-profile investigative hearings
into Wall Street practices that followed the crash, questioning many of the
book's banking figures directly about their conduct during the boom
\autocite{libraryjournal1929,dailyeconomy1929}. His hearings serve as the
narrative bridge between the crash and the reforms that followed it.

\end{description}

\section{Cast of Institutions}
\label{sec:cast-institutions}

\begin{description}

\item[National City Bank] A forerunner of Citibank and, under
\charname{Charles Mitchell}'s leadership, one of the most aggressive
commercial banks in extending credit for stock speculation and marketing
securities directly to retail customers \autocite{airmail1929,
supersummary1929}. Its practices put it repeatedly at odds with the
\charname{Federal Reserve}.

\item[J.P.\ Morgan \& Co.] The most powerful private bank on Wall Street, led
by \charname{Thomas Lamont} and \charname{J.P.\ Morgan Jr.}, whose partners
acted as informal coordinators of the banking community's response to market
stress \autocite{supersummary1929}. The firm's influence connects private
finance directly to the \charname{New York Stock Exchange} and to Washington.

\item[New York Stock Exchange (NYSE)] The physical and symbolic center of
American securities trading, where figures like \charname{Richard Whitney}
and \charname{Michael Meehan} operated, and whose self-regulated trading
floor practices are a recurring subject of the book
\autocite{supersummary1929}. It is the stage on which the boom's speculative
excesses and the crash itself play out.

\item[The Federal Reserve System] Comprising the Board in Washington and
regional banks such as the \charname{Federal Reserve Bank of New York} (led
by \charname{George Harrison}), the Fed grappled throughout the boom with
whether and how to restrain speculative lending by banks like
\charname{National City Bank} \autocite{supersummary1929}. Its internal
debates connect monetary policy to the conduct of individual bankers profiled
elsewhere in the book.

\item[General Motors] The automaker at the center of two major characters'
fortunes --- founder \charname{William C.\ Durant} and executive
\charname{John J.\ Raskob} --- General Motors' stock was itself one of the
era's most actively speculated-upon issues, tying industrial America
directly to the market's mania \autocite{supersummary1929}.

\item[The White House and U.S.\ Treasury] The executive branch under
\charname{Herbert Hoover} and Treasury Secretary \charname{Andrew Mellon},
and later under \charname{Franklin Roosevelt}, whose evolving posture toward
Wall Street --- from voluntary cooperation to active regulation --- structures
the political half of the book's narrative \autocite{supersummary1929,
csmonitor1929}.

\item[U.S.\ Senate Committee on Banking and Currency] The congressional
committee whose investigative hearings, led by chief counsel
\charname{Ferdinand Pecora} and shaped legislatively by Senator
\charname{Carter Glass}, became the public forum in which Wall Street's
boom-era practices were examined and which produced the era's landmark
banking reforms \autocite{libraryjournal1929,dailyeconomy1929}.

\end{description}

\section{Chapter-by-Chapter Narrative Outline}
\label{sec:narrative-outline}

\emph{This narrative outline follows SuperSummary's published chapter groupings
as a discursive, thematic read of the book's arc. The book's own chapter titles
and the front-matter cast list are reproduced directly, without this
intermediary, in Sections~\ref{sec:cast-companies} and~\ref{sec:chapter-titles}
below; see Section~\ref{sec:notes} for how the two are sourced differently. Each
entry below stays conceptual and omits how any individual character's story
resolves.}

\begin{enumerate}[leftmargin=1.6em, itemsep=0.6em]

\item \textbf{Prologue.} The book opens in medias res on the eve of the crash's worst
days in late October 1929, with \charname{Charles Mitchell} and \charname{National
City Bank} already under acute strain --- a framing device that sets up the question
the rest of the book answers: how the country arrived at this moment
\autocite{supersummaryPrologueCh8}.

\item \textbf{Part 1, Chapters 1--8 (February--spring 1929).} The narrative rewinds
to the height of the bull market. \charname{Charles Mitchell} and \charname{National
City Bank} keep credit flowing to speculators even as the \charname{Federal Reserve}
grows uneasy, and \charname{Herbert Hoover} takes office in March facing a market he
privately distrusts \autocite{wikipedia1929book,supersummaryPrologueCh8}.

\item \textbf{Part 1, Chapters 9--17 (spring--summer 1929).} The conflict between
Wall Street and Washington sharpens: \charname{William C.\ Durant} travels secretly
to Washington to warn \charname{Herbert Hoover} that the \charname{Federal Reserve}'s
tightening threatens the economy, while \charname{Carter Glass} publicly clashes with
\charname{Charles Mitchell} over the dangers of speculative bank credit, splitting
establishment opinion \autocite{supersummaryCh9to17}.

\item \textbf{Part 1, Chapters 18--25 (October 1929).} As confident public voices
insist stock prices have reached a permanently high plateau, \charname{Thomas Lamont}
and his \charname{J.P.\ Morgan \& Co.} partners privately debate how fragile the boom
has become. When panic selling breaks out on the \charname{New York Stock Exchange},
Lamont organizes the leading bankers into a rescue pool, with \charname{Richard
Whitney} acting as its public face on the trading floor, but the effort cannot hold
back the record-breaking collapse that follows within days
\autocite{supersummaryCh18to25}.

\item \textbf{Part 2, Preface--Chapter 33 (late 1929 -- early 1930s).} In the crash's
wake, \charname{Herbert Hoover}, \charname{Andrew Mellon}, and \charname{George
Harrison} confront a deepening economic contraction and a spreading wave of bank
failures, testing the limits of voluntary, business-led responses, while
\charname{Franklin D.\ Roosevelt} emerges as the political figure whose 1932 election
promises a different relationship between Washington and Wall Street
\autocite{supersummary1929,csmonitor1929}.

\item \textbf{Part 2, Chapter 34--Afterword (1933--1934 and beyond).}
\charname{Ferdinand Pecora}'s Senate hearings put \charname{Charles Mitchell} and
other boom-era figures under public questioning about their conduct, while
\charname{Carter Glass} and the \charname{Roosevelt} administration use the hearings'
momentum to reshape American banking law; the Afterword steps back to consider what
the era's crash and reforms mean for the financial crises that came after
\autocite{supersummaryCh34Afterword,axios1929}.

\end{enumerate}

\section{The Cast of Characters and the Companies They Kept}
\label{sec:cast-companies}

The book opens with its own roster of the people and institutions it follows,
organized by real-world affiliation rather than by narrative role. It is
reproduced in full below, following the book's own category headings and
grouping, as a quick-reference companion to the analytical sketches in
Sections~\ref{sec:cast-characters} and~\ref{sec:cast-institutions} above
\autocite{sorkinPreviewPDF}. Where the book itself gives a name with no further
description, none is added here.

\paragraph{Financial Institutions}

\subcat{The House of Morgan}
\begin{description}
\item[John Pierpont Morgan Sr.] Founder.
\item[J.\,P.\ ``Jack'' Morgan Jr.] Partner, son of John Pierpont Morgan.
\item[Thomas William Lamont] Partner.
\item[Thomas Stilwell Lamont] Partner, son of Thomas William Lamont.
\item[George Whitney] Partner.
\item[Henry Davison] Partner.
\item[Dwight Morrow] Former partner, ambassador to Mexico.
\item[Russell Cornell Leffingwell] Partner.
\item[Frank Bartow] Partner.
\item[Parker Gilbert] Partner and associate.
\item[Martin Egan] Publicist.
\end{description}

\subcat{National City Company and National City Bank}
\begin{description}
\item[Charles Edwin Mitchell] Chairman and CEO.
\item[Hugh Baker] Head of National City's stock-trading unit.
\item[Gordon Rentschler] Bank president.
\end{description}

\subcat{County Trust Company}
\begin{description}
\item[James Riordan] Founder and president.
\end{description}

\subcat{Chase National Bank}
\begin{description}
\item[Albert H.\ Wiggin] Chairman and CEO.
\end{description}

\subcat{First National Bank of New York}
\begin{description}
\item[George F.\ Baker Sr.] Chairman.
\item[George F.\ Baker Jr.] Vice chairman.
\end{description}

\subcat{Bankers Trust}
\begin{description}
\item[Seward Prosser] Chairman.
\end{description}

\subcat{Irving Trust}
\begin{description}
\item[Lewis Pierson] Chairman.
\end{description}

\subcat{Reichsbank}
\begin{description}
\item[Dr.\ Hjalmar Schacht] President.
\end{description}

\paragraph{Business Leaders}
\begin{description}
\item[Colonel Sosthenes Behn] Founder and chairman of International Telephone
and Telegraph.
\item[Louis-Joseph Chevrolet] Co-founder of Chevrolet Motor Company.
\item[Walter Percy Chrysler] Founder of Chrysler Corporation.
\item[Pierre Samuel du Pont] Board member of DuPont and General Motors.
\item[T.\ Coleman du Pont] President of E.\,I.\ du Pont de Nemours and
Company, two-term U.S.\ senator from Delaware.
\item[Henry Ford] Founder of Ford Motor Company.
\item[William Fox] Founder of Fox Film Corporation.
\item[John Jakob Raskob] Executive at DuPont and General Motors and chairman
of the Democratic National Committee (1928--1932).
\item[Colonel William P.\ Rend] President of the W.\,P.\ Rend Coal Company.
\item[David Sarnoff] Founding general manager of Radio Corporation of America
(RCA).
\item[Charles M.\ Schwab] President of Bethlehem Steel.
\item[Alfred Sloan] President of General Motors.
\item[Owen D.\ Young] President and chairman of General Electric, co-founder
of Radio Corporation of America (RCA).
\end{description}

\paragraph{The New York Stock Exchange}
\begin{description}
\item[E.\,H.\,H.\ Simmons] President (1924--1930).
\item[Richard Whitney] Vice president, broker for J.P.\ Morgan \& Co.
\item[William Crawford] Superintendent.
\item[Michael Meehan] RCA, or ``Radio,'' specialist and pool operator.
\item[General Oliver Bridgeman] U.S.\ Steel specialist.
\item[Ben Smith] Pool operator.
\item[Tom Bragg] Pool operator.
\item[Michael Levine] Overseer of Wall Street's army of messenger boys.
\end{description}

\paragraph{Stockbrokers}
\begin{description}
\item[Charles Merrill] Merrill Lynch co-founder.
\item[William Van Antwerp] Partner at E.\,F.\ Hutton and Co.
\end{description}

\paragraph{The Speculators}
\begin{description}
\item[Evangeline Adams] Astrologist known as the ``stock market's seer.''
\item[Bernard Baruch]
\item[Pat Bologna] Wall Street bootblack with a reputation as a tipster.
\item[Arthur Cutten]
\item[William Crapo Durant]
\item[Clarence Hatry] British financier.
\item[Joseph Kennedy]
\item[Jesse Livermore]
\item[Groucho Marx] Actor.
\item[Oris and Mantis Van Sweringen]
\end{description}

\paragraph{New York}
\begin{description}
\item[Jimmy Walker] Mayor of New York City (1926--1932).
\item[Fiorello La Guardia] Mayor of New York City (1934--1946).
\item[Al Smith] Governor of New York (1919--1920, 1923--1928), Democratic
candidate for president in 1928.
\end{description}

\paragraph{United Kingdom}
\begin{description}
\item[Winston Churchill] British member of Parliament, former chancellor of
the exchequer, future prime minister (1940--1945, 1951--1955).
\item[Randolph Churchill] Son of Winston Churchill.
\item[Ramsay MacDonald] British Labour Party leader, foreign secretary, prime
minister (1924, 1929--1935).
\end{description}

\paragraph{U.S.\ Government}
\begin{description}
\item[William Howard Taft] Chief justice of the U.S.\ Supreme Court
(1921--1930), U.S.\ president (1909--1913).
\item[Warren G.\ Harding] U.S.\ president (1921--1923).
\item[Calvin Coolidge] U.S.\ president (1923--1929).
\item[Herbert Clark Hoover] U.S.\ president (1929--1933).
\item[Franklin Delano Roosevelt] Governor of New York (1929--1933), U.S.\
president (1933--1945).
\item[Andrew William Mellon] American banker and businessman, treasury
secretary (1921--1932).
\item[William Woodin] Industrialist, treasury secretary (1933).
\item[Senator Carter Glass] Democrat from Virginia, co-founder of the Federal
Reserve Act of 1913, co-sponsor of the Glass--Steagall Act of 1932.
\item[Senator Hiram Johnson] Republican from California.
\item[Senator Huey Long] Nicknamed ``the Kingfish,'' Democrat from Louisiana.
\item[Senator Kenneth McKellar] Democrat from Tennessee.
\item[Senator George Moses] Republican from New Hampshire.
\item[Senator George Norris] Republican from Nebraska.
\item[Senator Robert L.\ Owen] Democrat from Oklahoma, co-sponsor with Carter
Glass of the Federal Reserve Act.
\item[Senator Arthur Robinson] Republican from Indiana.
\item[Senator Reed Owen Smoot] Republican from Utah, co-sponsor of the 1930
Smoot--Hawley Tariff Act.
\item[Representative Henry B.\ Steagall] Democrat from Alabama, co-sponsor of
the Glass--Steagall Act of 1932.
\item[Learned Hand] Federal judge on the U.S.\ Court of Appeals for the
Second Circuit.
\item[Harold R.\ Medina] Federal judge on the U.S.\ District Court for the
Southern District of New York.
\end{description}

\paragraph{The Federal Reserve}
\begin{description}
\item[Marriner S.\ Eccles] Chairman of the Federal Reserve Board
(1934--1948).
\item[George L.\ Harrison] Governor and then president of the Federal
Reserve Bank of New York (1928--1940).
\item[Adolph C.\ Miller] Member of the Federal Reserve Board (1914--1936).
\item[Benjamin Strong Jr.] Governor of the Federal Reserve Bank of New York
(1914--1928).
\item[Charles Sumner Hamlin] First chairman of the Federal Reserve Board
(1914--1916), remained a member of the board until 1936.
\item[Paul M.\ Warburg] Head of the International Acceptance Bank, involved
in creating the Federal Reserve, on the first Washington board of governors.
\end{description}

\paragraph{The Investigators}
\begin{description}
\item[William O.\ Douglas] Chairman of the Securities and Exchange
Commission (SEC), future associate justice of the U.S.\ Supreme Court.
\item[Ferdinand Pecora] New York prosecutor known as the ``Hellhound of Wall
Street'' for leading an investigation of the city's bankers in the wake of
the 1929 crash.
\item[Ars\`ene P.\ Pujo] Louisiana congressman and head of the Committee on
Banking and Currency, chaired the ``Pujo Committee'' investigation into the
financial industry in 1912.
\end{description}

\paragraph{The Economists}
\begin{description}
\item[Roger Babson] Statistician and economist who predicted the crash.
\item[Irving Fisher] Yale professor, one of the nation's leading economists.
\item[John Kenneth Galbraith] Harvard professor of economics.
\item[John Maynard Keynes] British economist.
\item[Joseph Stagg Lawrence] Princeton economist.
\item[Royal Meeker] Pro--stock market economist.
\end{description}

\paragraph{The Journalists}
\begin{description}
\item[Claud Cockburn] British writer for \emph{The Times} of London.
\item[Samuel Crowther] Author of ``Everybody Ought to Be Rich'' for
\emph{Ladies' Home Journal}.
\item[William Floyd] Author of \emph{People vs.\ Wall Street}.
\item[Major Robert H.\ Glass] Newspaperman, father of Carter Glass.
\item[William Randolph Hearst] Publisher of Hearst Newspapers.
\item[Matthew Josephson] Writer who had worked in the financial industry.
\item[Edwin Lef\`evre] Author of \emph{Reminiscences of a Stock Operator}.
\item[Walter Lippmann] Writer, reporter, and political commentator.
\item[Charley Michelson] Hearst newspaper writer turned Democratic publicist.
\item[Alexander D.\ Noyes] Financial editor of \emph{The New York Times}.
\item[Drew Pearson] Syndicated columnist of \emph{Washington Merry-Go-Round}.
\end{description}

\section{Chapter Titles}
\label{sec:chapter-titles}

The book's own table of contents lists a Prologue, forty-two dated chapters
across two parts, an Epilogue, and an Afterword \autocite{vitalsource1929}. Each
entry below gives the chapter's actual title and one sentence of independently
researched historical context for its date --- not a summary of the chapter's
own contents, which were not available to this document's author beyond the
Prologue excerpted in the publisher preview.

\begin{description}
\item[Prologue] Set on the evening of Monday, October 28, 1929, hours after the
market's worst single-session decline to that point (a 13 percent drop
remembered as Black Monday, the day before the still-steeper 12 percent decline
of Black Tuesday), the Prologue opens after the crash the rest of the book
builds toward \autocite{sorkinPreviewPDF}.
\end{description}

\subsection*{Part I}

\begin{description}

\item[Chapter One: February 1, 1929] By early 1929, brokers' loans financing
stock speculation had roughly tripled over two years, prompting a worried
Federal Reserve Board to begin closed-door discussions on restraining bank
lending into the market \autocite{fedhistorycrash}.

\item[Chapter Two: February 14, 1929] Through mid-February 1929 the Federal
Reserve Board weighed, and repeatedly held back from, a direct public warning
against speculative lending, even as the market continued climbing
\autocite{fedhistorycrash}.

\item[Chapter Three: February 16, 1929] The Federal Reserve Bank of New York,
under governor George Harrison, had begun pressing Washington to let it raise
its discount rate to cool speculation, a request the Board in Washington
continued to resist through the winter of 1929 \autocite{fedhistorycrash}.

\item[Chapter Four: March 4, 1929] Herbert Hoover was inaugurated as the
thirty-first president of the United States on March 4, 1929, inheriting a
stock market boom that had continued unbroken through the transition from
Coolidge's presidency \autocite{crashwiki}.

\item[Chapter Five: March 5, 1929] In his first days in office, Hoover
inherited an unresolved fight inside his own Federal Reserve over how
aggressively to restrain the boom he now presided over \autocite{fedhistorycrash}.

\item[Chapter Six: March 26, 1929] On March 26, 1929, call-money rates spiked
as high as 20 percent, triggering a one-day panic that Charles Mitchell halted
by pledging \$25 million of National City Bank's own funds to the loan market
--- a public rebuke to the Federal Reserve that drew the immediate anger of
Senator Carter Glass \autocite{mitchellmarch1929}.

\item[Chapter Seven: March 29, 1929] In the days after Mitchell's
intervention, call money eased from 20 percent back toward single digits and
the immediate panic passed, though the underlying tension between Wall Street
and the Fed did not \autocite{mitchellmarch1929}.

\item[Chapter Eight: April 5, 1929] Through early April 1929 the bull market
resumed its climb, with both stock prices and margin debt setting fresh
records even as the Federal Reserve's internal fight over discount rates
remained unresolved \autocite{fedhistorycrash}.

\item[Chapter Nine: April 8, 1929] The New York Fed continued to press through
the spring of 1929 for authority to raise its discount rate, a request the
Board in Washington would not grant until August \autocite{fedhistorycrash}.

\item[Chapter Ten: April 12, 1929] General Motors and Radio Corporation of
America remained among the most actively traded and speculated-upon stocks of
the spring, emblematic of the boom's favorite issues \autocite{crashwiki}.

\item[Chapter Eleven: April 14, 1929] By mid-April 1929, total brokers' loans
--- the credit fueling margin speculation --- had climbed past levels that
even boosters of the market found difficult to fully explain
\autocite{fedhistorycrash}.

\item[Chapter Twelve: May 7, 1929] Through May 1929 the standoff between the
Federal Reserve Board and the New York Fed over how to cool the market
continued without resolution, even as prices kept climbing
\autocite{fedhistorycrash}.

\item[Chapter Thirteen: June 4, 1929] By early June 1929, the market's advance
had continued for so long that some economists, including Yale's Irving
Fisher, had begun publicly arguing stock prices reflected a durable new
economic reality rather than speculation \autocite{crashwiki}.

\item[Chapter Fourteen: June 29, 1929] As the first half of 1929 closed, the
market had already climbed roughly 25 percent for the year, with little sign
yet of the reversal to come \autocite{crashwiki}.

\item[Chapter Fifteen: September 2, 1929] September 2, 1929, was the eve of
the market's all-time pre-crash high: the Dow Jones Industrial Average closed
at 381.17 the following day, September 3 \autocite{crashwiki}.

\item[Chapter Sixteen: October 2, 1929] By early October 1929, the market's
advance had stalled and volatility had increased, though few observers yet
recognized it as the beginning of the crash \autocite{crashwiki}.

\item[Chapter Seventeen: October 6, 1929] Selling pressure continued through
the first week of October 1929, with the market down significantly from its
September peak \autocite{crashwiki}.

\item[Chapter Eighteen: October 10, 1929] On October 10, 1929, J.P.\ Morgan
partner Thomas Lamont hosted British prime minister Ramsay MacDonald for
dinner in New York while privately debating with partner Russell Leffingwell
how fragile the boom had become \autocite{supersummaryCh18to25}.

\item[Chapter Nineteen: October 24, 1929] October 24, 1929 --- Black Thursday
--- brought the worst single-session panic the market had yet seen, prompting
Lamont to convene the ``Big Six'' bankers into a pool of more than \$200
million to support prices \autocite{supersummaryCh18to25}.

\item[Chapter Twenty: October 27, 1929] October 27, 1929, was the Sunday
between Black Thursday and the still-steeper declines of Black Monday and
Black Tuesday that followed on October 28 and 29 \autocite{crashwiki}.

\item[Chapter Twenty-One: November 6, 1929] In the week after Black Tuesday,
the market continued grinding lower toward the bottom it would reach in
mid-November \autocite{crashwiki}.

\item[Chapter Twenty-Two: November 8, 1929] Selling continued into the second
week of November 1929, with the Dow still well short of the low it would set
on November 13 \autocite{crashwiki}.

\item[Chapter Twenty-Three: November 13, 1929] The Dow Jones Industrial
Average closed at 198.60 on November 13, 1929 --- its lowest point of the
crash, down roughly 48 percent from its September peak \autocite{crashwiki}.

\item[Chapter Twenty-Four: December 19, 1929] By mid-December 1929, the
National Business Survey Conference that Hoover had tasked with enforcing the
wage- and investment-holding pledges he secured from industry leaders that
November was just getting underway \autocite{hooverconferences1929}.

\item[Chapter Twenty-Five: December 21, 1929] As 1929 drew to a close, the
market had stabilized well below its autumn highs, though the deeper economic
contraction to come was not yet apparent to most observers \autocite{crashwiki}.

\end{description}

\subsection*{Part II}

\begin{description}

\item[Chapter Twenty-Six: September 30, 1930] Nearly a year after the crash,
early hopes of a quick recovery had given way to a spreading wave of regional
bank failures, the first of what the Federal Reserve later identified as the
banking panics of 1930--31 \autocite{fedhistorypanics3031}.

\item[Chapter Twenty-Seven: October 28, 1930] Bank failures continued
accelerating through the fall of 1930, setting the stage for the collapse
that December of New York's Bank of United States \autocite{bankofusfail}.

\item[Chapter Twenty-Eight: November 5, 1930] Deposit runs continued at banks
nationwide in the weeks before the Bank of United States failed on December
11, 1930, then the largest bank failure in U.S.\ history \autocite{bankofusfail}.

\item[Chapter Twenty-Nine: February 2, 1931] A relative lull in bank failures
in early 1931 proved temporary; a fresh, still larger wave of the banking
panic was still to come later that year \autocite{fedhistorypanics3031}.

\item[Chapter Thirty: February 18, 1932] This chapter falls between two
Hoover-era interventions: his creation of the Reconstruction Finance
Corporation to lend directly to troubled banks on January 22, 1932, and his
signing, on February 27, of the Banking Act of 1932
\autocite{fedhistorybankingact1932,rfcwiki}.

\item[Chapter Thirty-One: November 8, 1932] Franklin D.\ Roosevelt was elected
president in a landslide over the incumbent Hoover on November 8, 1932, a
verdict widely read as a repudiation of the administration's handling of the
Depression \autocite{crashwiki}.

\item[Chapter Thirty-Two: February 18, 1933] Four days earlier, on February
14, 1933, Michigan's governor had declared the nation's first statewide bank
holiday, closing every bank in the state amid the collapse of Detroit's Union
Guardian Trust \autocite{michiganholiday1933}.

\item[Chapter Thirty-Three: February 21, 1933] In the days that followed
Michigan's lead, additional states began declaring their own bank holidays as
confidence in the banking system eroded further \autocite{michiganholiday1933}.

\item[Chapter Thirty-Four: February 22, 1933] By late February 1933, the wave
of state bank holidays was continuing to spread in the final days of Hoover's
presidency \autocite{michiganholiday1933}.

\item[Chapter Thirty-Five: March 3, 1933] On the eve of Roosevelt's
inauguration, New York and Illinois joined dozens of other states in
declaring bank holidays, effectively shutting down the national banking
system \autocite{emergencybankingact1933wiki}.

\item[Chapter Thirty-Six: March 7, 1933] Roosevelt was sworn in on March 4,
1933, and declared a nationwide bank holiday on March 6, closing every bank
in the country while Congress prepared emergency legislation
\autocite{emergencybankingact1933wiki}.

\item[Chapter Thirty-Seven: March 21, 1933] By late March 1933, Congress had
passed the Emergency Banking Act (March 9) and Roosevelt had delivered his
first fireside chat (March 12), reassuring depositors it was safer to keep
their money in a reopened bank than under a mattress
\autocite{emergencybankingact1933wiki}.

\item[Chapter Thirty-Eight: May 16, 1933] Through May 1933, Congress worked
toward the Securities Act, which would require public companies to disclose
financial information to investors for the first time
\autocite{securitiesact1933wiki}.

\item[Chapter Thirty-Nine: May 22, 1933] Days before the Securities Act's
signing, congressional negotiations over its final disclosure requirements
continued \autocite{securitiesact1933wiki}.

\item[Chapter Forty: May 23, 1933] With the Securities Act's passage
imminent, its requirement that new securities be registered with federal
regulators marked the first federal effort to police what companies told
investors \autocite{securitiesact1933wiki}.

\item[Chapter Forty-One: June 16, 1933] On June 16, 1933, Roosevelt signed
the Banking Act of 1933 --- popularly known as Glass-Steagall --- separating
commercial from investment banking and creating the Federal Deposit Insurance
Corporation \autocite{glasssteagall1933wiki}.

\item[Chapter Forty-Two: June 21, 1933] Days after Glass-Steagall's signing,
regulators began the work of implementing its new deposit-insurance and
bank-separation rules \autocite{glasssteagall1933wiki}.

\end{description}

\begin{description}
\item[Epilogue] Per the book's table of contents, the Epilogue follows the
final dated chapter as a distinct, reflective closing section
\autocite{vitalsource1929}.
\item[Afterword] The Afterword closes the volume proper, ahead of the
photographs, acknowledgments, and source notes that follow in the back matter
\autocite{vitalsource1929}.
\end{description}

\section{Notes on This Document}
\label{sec:notes}

This companion is a single LaTeX source (\texttt{sorkin-1929-reader-companion.tex})
built with \texttt{pdflatex} and \texttt{biber} against a BibTeX bibliography
(\texttt{references.bib}); citations and the numbered References list use
\texttt{biblatex}'s \texttt{numeric} style, ordered by first appearance. A
companion EPUB edition is produced by the same source tree's \texttt{build.sh}
script: \texttt{tex4ebook} is used directly when installed, but since it ships
only from CTAN --- unreachable under this authoring environment's network
policy, and not packaged for Ubuntu --- the script falls back to the engine
\texttt{tex4ebook} itself wraps: \texttt{htlatex} (tex4ht) \textrightarrow\
\texttt{biber} \textrightarrow\ \texttt{htlatex} \textrightarrow\ \texttt{pandoc}
(HTML to EPUB3). That fallback output validates cleanly under EPUBCheck.

This document's sourcing has two distinct tiers, kept separate rather than
reconciled, so a reader can see which claims rest on the book itself and which
rest on how others have described it. Sections~\ref{sec:cast-characters},
\ref{sec:cast-institutions}, and~\ref{sec:narrative-outline} were built from
publicly available secondary sources --- publisher copy, book reviews, press
interviews, and reader/study-guide summaries --- because neither the book's
text nor its full table of contents was available when they were first
drafted; the narrative outline in particular follows SuperSummary's published
chapter groupings (Prologue; Part 1, Chapters 1--8, 9--17, and 18--25; Part 2,
Preface--Chapter 33; and Chapter 34--Afterword) as the closest verifiable proxy
for a table of contents once ursummary.com's own page proved unreachable from
this document's authoring environment \autocite{ursummaryChapters1929}.
Sections~\ref{sec:cast-companies} and~\ref{sec:chapter-titles} were added in a
later revision once the book's own front-matter cast list and complete
table of contents became available via a publisher preview PDF and a
VitalSource export the reader supplied directly, and reproduce that primary
material rather than a secondary summary of it
\autocite{sorkinPreviewPDF,vitalsource1929}. Section~\ref{sec:about} was
revised in that same pass and mixes both tiers: its biographical framing of
Sorkin still rests on secondary sources, but its account of the book's stated
influences and research process is drawn directly from his author's note in
the same preview PDF. The one-sentence notes attached
to each chapter title in Section~\ref{sec:chapter-titles} describe the
real-world historical backdrop of that chapter's date, independently
researched, rather than the chapter's own contents, which were not available
to this document's author beyond the Prologue.

Throughout, character and institution sketches describe who people were and
how they connect to one another, deliberately omitting how each person's
individual story concludes. This document is not authorized, reviewed, or
endorsed by the author or publisher, and no part of the book's own text is
reproduced here beyond the front-matter cast list and chapter titles
explicitly requested.

\printbibliography[title={References}]

\end{document}
